\documentclass[a4paper,11pt]{article}
        \usepackage[top=15mm,bottom=18mm,left=16mm,right=16mm]{geometry}
        \usepackage{fontspec}
        \usepackage{xeCJK}
        \setmainfont{Times New Roman}
        \setCJKmainfont{Yu Gothic}
        \setmonofont{Consolas}
        \setCJKmonofont{Yu Gothic}
        \usepackage{booktabs}
        \usepackage{longtable}
        \usepackage{tabularx}
        \usepackage{array}
        \usepackage{enumitem}
        \usepackage{hyperref}
        \usepackage{listings}
        \usepackage{xcolor}
        \usepackage{amsmath}
        \usepackage{amssymb}
        \hypersetup{
          colorlinks=true,
          linkcolor=blue,
          urlcolor=blue,
          pdftitle={live 風補正ノード},
          pdfauthor={Codex}
        }
        \lstset{
          basicstyle=\ttfamily\small,
          breaklines=true,
          columns=fullflexible,
          keepspaces=true,
          frame=single,
          framerule=0.2pt,
          rulecolor=\color{black},
          showstringspaces=false
        }
        \setlength{\parskip}{0.42em}
        \setlength{\parindent}{1em}
        \setlength{\emergencystretch}{3em}
        \renewcommand{\arraystretch}{1.15}
        \title{23. live 風補正ノード}
        \author{solar\_ws0129-main}
        \date{2026-07-29}
        \begin{document}
        \maketitle
        \tableofcontents
        \clearpage

        \section{対象}
        \begin{tabularx}{\linewidth}{>{\raggedright\arraybackslash}p{0.18\linewidth} X}
        \toprule
        項目 & 内容 \\
        \midrule
        ファイル & \path|mpc_solarcar/wind_correction_node.py| \\
        ソースSHA-256 & \path|065128cbb44604d15e3d04030f00c23607abfa9faf13ebb66e08835897217593| \\
        種別 & Python \\
        区分 & runtime node \\
        役割 & 観測 headwind と現在距離から raw forecast CSV を補正し、corrected forecast CSV を書き出す。 \\
        起動文脈 & live\_wifi で planner 入力の風を上書きする前処理。 \\
        \bottomrule
        \end{tabularx}

        \section{このファイルを読む理由}
        観測 headwind と現在距離から raw forecast CSV を補正し、corrected forecast CSV を書き出す。

        \begin{itemize}[leftmargin=1.6em]
        \item planner は /weather/headwind\_corrected\_ms を直接読むのではない。
\item corrected CSV を mpc\_node が再読込して効かせる。
        \end{itemize}

        \section{呼び出し関係}
        \subsection{どこから呼ばれるか}
        \begin{itemize}[leftmargin=1.6em]
        \item \path|launch/solar_race_live_wifi.launch.py|
        \end{itemize}

        \subsection{次に読むと理解が進むファイル}
        \begin{itemize}[leftmargin=1.6em]
        \item \path|mpc_solarcar/mpc_node.py|
        \end{itemize}

        \subsection{同一リポジトリ内の主要依存}
        \begin{itemize}[leftmargin=1.6em]
        \item 同一リポジトリ内の明示 import は少ない。
        \end{itemize}

        \section{ソース冒頭の把握}
        モジュール docstring または先頭意図の要約:

        モジュール docstring は明示されていない。

        \subsection{代表コード断片}
        \begin{lstlisting}[language=Python]
return int(pd.Timestamp(value).value)


class WindCorrectionNode(Node):
    def __init__(self) -> None:
        super().__init__("wind_correction_node")
        self.declare_parameter("forecast_csv", "")
        self.declare_parameter("corrected_forecast_csv", "outputs/runtime/live_forecast_corrected.csv")
        self.declare_parameter("forecast_time_tz", "Australia/Darwin")
        self.declare_parameter("publish_period_sec", 30.0)
        self.declare_parameter("measurement_sigma_ms", 1.0)
        self.declare_parameter("correlation_distance_km", 300.0)
        self.declare_parameter("fallback_correlation_time_h", 3.0)
        self.declare_parameter("forecast_sigma0_ms", 1.5)
        self.declare_parameter("forecast_variance_growth_per_hour", 0.05)
        self.declare_parameter("planning_quantile", 0.5)
        self.declare_parameter("confidence_z", 1.96)
        self.declare_parameter("min_sigma_ms", 0.2)
        self.declare_parameter("preferred_source", "auto")
        self.declare_parameter("use_exp_distance_decay", True)

        self.forecast_csv = Path(str(self.get_parameter("forecast_csv").value or "")).expanduser()
\end{lstlisting}


        \section{主要構造の読み取り}
        主要クラスは WindCorrectionNode。 主要関数は main。 ROS パラメータ宣言は 14 件。 ROS I/O は publisher 2 件、subscription 2 件。

        \subsection{主要クラス・関数}
            \begin{longtable}{p{0.07\linewidth} p{0.16\linewidth} p{0.25\linewidth} p{0.40\linewidth}}
            \toprule
            行 & 種別 & 名前 & 補足 \\
            \midrule
            20 & class & \path|WindCorrectionNode| & bases: Node \\
16 & function & \path|_timestamp_ns| & args: value \\
21 & function & \path|__init__| & args: self \\
64 & function & \path|_on_headwind| & args: self, msg \\
67 & function & \path|_on_s| & args: self, msg \\
70 & function & \path|_load_base| & args: self \\
91 & function & \path|_interp_headwind_now| & args: self, now\_utc \\
105 & function & \path|_step| & args: self \\
154 & function & \path|main| &  \\
            \bottomrule
            \end{longtable}

        \subsection{ファイルを上から読んだときの定義順}
            \begin{longtable}{p{0.07\linewidth} p{0.84\linewidth}}
            \toprule
            行 & 内容 \\
            \midrule
            16 & 関数 \_timestamp\_ns を定義する。 \\
20 & クラス WindCorrectionNode を定義する。 \\
154 & 関数 main を定義する。 \\
            \bottomrule
            \end{longtable}

        \subsection{import 群の詳細}
            \begin{longtable}{p{0.07\linewidth} p{0.33\linewidth} p{0.50\linewidth}}
            \toprule
            行 & import 文 & このファイルで読む理由 \\
            \midrule
            1 & \path|from __future__ import annotations| & 型ヒントの遅延評価を有効にし、自己参照型や forward reference を安全に扱うため。 このファイル内での使用位置は少ないか、間接利用である。 \\
3 & \path|import math| & Python標準のスカラー数値関数と定数を使うため。実際に参照する名前と行は使用位置欄で確認する。 このファイル内での主な使用位置は L52, L53, L111, L121, L122, L130, L131, L135。 \\
4 & \path|import os| & パス、環境変数、プロセス外部状態を扱うため。 このファイル内での主な使用位置は L72。 \\
5 & \path|from datetime import datetime, timezone| & UTC 時刻や相対時間を扱うため。 このファイル内での主な使用位置は L91, L109, L128。 \\
6 & \path|from pathlib import Path| & ファイルやディレクトリを安全に扱うため。 このファイル内での主な使用位置は L38, L39。 \\
7 & \path|from statistics import NormalDist| & statistics から NormalDist を読み込み、このファイルの処理を組み立てるため。 このファイル内での主な使用位置は L50。 \\
9 & \path|import pandas as pd| & CSV や時刻列の読込・整形・再サンプリングに使うため。 このファイル内での主な使用位置は L17, L77, L79, L96, L119, L128, L140, L147。 \\
10 & \path|import rclpy| & ROS 2 Python ノードとして起動・spin するため。 このファイル内での主な使用位置は L155, L157, L159。 \\
11 & \path|from rclpy.node import Node| & ROS 2 ノード本体の基底クラスとして使うため。 このファイル内での主な使用位置は L20。 \\
13 & \path|from std_msgs.msg import Float32, String| & Float32、Bool、Stringなどの標準ROS messageで値をpublishまたはsubscribeするため。 このファイル内での主な使用位置は L57, L58, L59, L60, L64, L67, L115, L151。 \\
            \bottomrule
            \end{longtable}

        \section{このファイルを読む前に必要な基礎知識}
        次の章は、構文やROS用語を既知と仮定しないための説明である。

        \subsection{プログラム、プロセス、メモリ、オブジェクトを区別する}
ソースファイルはディスク上の文字列であり、それ自体は走っていない。Python実行可能プログラムがソースを読み、OSがその実行を一つのプロセスとして管理する。プロセスには仮想メモリ、開いているファイル、スレッド、終了コードなどが対応する。


\begin{lstlisting}
ソースファイル -> Pythonインタプリタが読み込む -> OS上のプロセス -> Pythonオブジェクトをメモリに生成 -> 関数やcallbackを実行
\end{lstlisting}


「メモリ上に生成する」とは、実行中プロセスが使う記憶領域に、その値の型、属性、参照関係を表す実体を用意することである。変数は実体そのものというより、そのオブジェクトを指す名前として理解するとPythonの代入が読みやすい。


\begin{lstlisting}[language=python]
node = MPCNode()
alias = node
# nodeとaliasは同じオブジェクトを参照する。
\end{lstlisting}


一つのプロセスには複数スレッドを持たせられる。スレッドは同じプロセスのメモリを共有するため、受信callbackと最適化callbackが同じself.zを書き換える場合は実行順と排他を検討する必要がある。

\paragraph{根拠資料}
\begin{itemize}[leftmargin=1.6em]
\item \href{https://docs.python.org/3/tutorial/classes.html}{Python公式チュートリアル: Classes}
\end{itemize}

\subsection{名前、代入、型、参照}
Pythonの代入文は、右辺を先に評価し、その結果のオブジェクトへ左辺の名前を結び付ける。`x = f()`では、まずfを呼び、その戻り値が得られてからxが更新される。


`float(...)`、`int(...)`、`bool(...)`は型名を呼び出して値を変換する。外部CSV、YAML、ROS parameterから来た値は型が期待どおりとは限らないため、このリポジトリでは明示変換が多い。


`None`は値が無いことを示す単一のオブジェクト、`math.nan`は浮動小数点値ではあるが有効な数値ではないことを示す。両者は用途が異なるため、`is None`と`math.isfinite`を使い分ける。

\paragraph{根拠資料}
\begin{itemize}[leftmargin=1.6em]
\item \href{https://docs.python.org/3/tutorial/classes.html}{Python公式チュートリアル: Classes}
\end{itemize}

\subsection{関数、引数、戻り値、スコープ}
`def`は関数本体をその場で実行する文ではない。関数オブジェクトを作り、その名前を現在の名前空間へ登録する。`f`は関数そのもの、`f()`はその関数を呼んだ結果である。


仮引数は関数定義側の受け取り名、実引数は呼出側から渡す値である。位置引数、キーワード引数、既定値付き引数、`*args`、`**kwargs`は渡し方が異なる。


\begin{lstlisting}[language=python]
def clip_speed(value, minimum=0.0, maximum=110.0):
    return min(maximum, max(minimum, value))

v = clip_speed(120.0, maximum=90.0)
\end{lstlisting}


関数内で作った通常の名前はローカルスコープに属する。メソッドが`self.z`を更新すればオブジェクトの状態が残るが、単に`z`を更新しただけならその呼出しのローカル値である。

\paragraph{根拠資料}
\begin{itemize}[leftmargin=1.6em]
\item \href{https://docs.python.org/3/tutorial/classes.html}{Python公式チュートリアル: Classes}
\end{itemize}

\subsection{module、package、import、console\_scripts}
Pythonファイルはmoduleとして読み込める。複数moduleをディレクトリにまとめたものがpackageである。`from .model import SolarCarModel`の先頭の点は、現在と同じpackage内のmodel moduleを指す相対importである。


import時にはファイルのトップレベル文が上から一度実行される。`def`や`class`は関数・クラスを登録するが、その本体の通常処理は呼び出すまで走らない。


setuptoolsの`console\_scripts`は、端末で使う実行可能名と`package.module:function`を対応付けるインストール時メタデータである。生成される実行用ラッパーはmoduleをimportし、指定関数を呼び、戻り値を終了コードとして扱う小さな入口である。


\begin{lstlisting}
ROS launchのexecutable名 -> install済みconsole script -> mpc_solarcar.mpc_nodeをimport -> main()を呼ぶ
\end{lstlisting}

\paragraph{根拠資料}
\begin{itemize}[leftmargin=1.6em]
\item \href{https://setuptools.pypa.io/en/latest/userguide/entry_point.html}{setuptools公式: Entry Points / Console Scripts}
\item \href{https://docs.python.org/3/tutorial/classes.html}{Python公式チュートリアル: Classes}
\end{itemize}

\subsection{例外、try/finally、with、資源解放}
例外は通常の戻り値とは別経路で異常を呼出元へ伝える。`try/except`は想定した異常を処理し、`finally`は成功・失敗にかかわらず後始末を行う。


`with open(...) as f:`はcontext managerを使い、ブロックを出るとファイルを閉じる。CSVやログの破損を避けるため、開いた資源の所有者と閉じる場所を明確にする。


`except Exception: pass`はノードを止めない利点がある一方、入力異常を隠して原因追跡を難しくする。安全に関係する値では、少なくとも頻度制限付き警告、異常カウンタ、fallback状態のpublishを検討する。



\subsection{class、独自クラス、インスタンス、self、継承}
クラスは、保持するデータと、そのデータを扱う関数を一つの型としてまとめる設計単位である。「独自クラス」とはPythonやROS 2が最初から用意した型ではなく、このプロジェクトが目的に合わせて定義した新しい型を指す。


`class MPCNode(Node)`は、rclpyのNodeを基底クラスとして継承し、Nodeが持つpublisher、subscription、timer、parameterなどの機能にMPC固有機能を追加する。丸括弧内は関数引数ではなく基底クラス指定である。


`MPCNode()`はクラスオブジェクトを呼び出して新しいインスタンスを作る。Pythonは新しい実体を用意し、その実体を第1引数selfとして`\_\_init\_\_`へ渡す。`self`は予約語ではなく慣習的な名前だが、この慣習を守ることでコードの意味が共有される。


\begin{lstlisting}[language=python]
class Counter:
    def __init__(self):
        self.value = 0

    def increment(self):
        self.value += 1

counter = Counter()
counter.increment()
\end{lstlisting}


`super().\_\_init\_\_(...)`は継承元の初期化処理を呼ぶ。MPCNodeではこれを省くとROSノードとして必要な内部実体が作られず、`create\_publisher`などを正しく使えない。

\paragraph{根拠資料}
\begin{itemize}[leftmargin=1.6em]
\item \href{https://docs.python.org/3/tutorial/classes.html}{Python公式チュートリアル: Classes}
\item \href{https://docs.ros.org/en/humble/Tutorials/Beginner-CLI-Tools/Understanding-ROS2-Nodes/Understanding-ROS2-Nodes.html}{ROS 2 Humble公式: Understanding nodes}
\end{itemize}

\subsection{先頭アンダースコア、\_\_init\_\_、名前付けの作法}
`self.\_step\_solar`の先頭1個のアンダースコアは「クラス内部で使う実装詳細」という慣習であり、アクセスを強制禁止する機構ではない。外部コードから呼べるが、公開APIとして依存しない意思表示である。


`\_\_init\_\_`の前後2個のアンダースコアはPythonが定めた特殊メソッド名である。クラス生成時に自動で呼ばれる。任意の名前に前後2個を付けて独自仕様を作ることは避ける。


`\_`で始まるローカル変数は、未使用または外部へ見せない意図を示す場合がある。例えば`\_base\_csv`は戻り値として受け取る必要はあるが、その後の処理では使わないことを読者へ示す。

\paragraph{根拠資料}
\begin{itemize}[leftmargin=1.6em]
\item \href{https://docs.python.org/3/tutorial/classes.html}{Python公式チュートリアル: Classes}
\end{itemize}

\subsection{list、dict、deque、NumPy配列、pandas DataFrame}
listは順序付き可変列、tupleは順序付きで通常変更しない列、dictはキーから値を引く対応表、dequeは両端追加・削除が効率的なキューである。どの構造を選ぶかはアクセス方法と更新方法で決まる。


NumPy配列は同種数値を連続的に扱い、要素ごとの演算、clip、補間、線形代数を簡潔に書く。shapeは各次元の要素数であり、速度系列なら通常`(N,)`、候補集団なら`(population, N)`となる。


pandas DataFrameは列名を持つ表である。CSV読込後は列型、欠損、単位、timezone、並び順、重複を明示的に処理しなければ、数値計算が動いても意味が誤る。



\subsection{rclpy、rcl、rmw、DDS/RTPSの層}
rclpyはPython利用者向けROS 2 client libraryである。利用者のNode、publisher、subscriptionなどをPython APIとして提供し、その下で共通C層のrclを利用する。


\begin{lstlisting}
本プロジェクトのPythonコード -> rclpy -> rcl -> rmw -> DDS/RTPS実装 -> ネットワークまたは同一PC内通信
\end{lstlisting}


rclは言語に依存しない共通ROS機能を提供するC API、rmwはROS 2と具体的middleware実装の境界である。DDS/RTPS側が探索、serialize、publish/subscribe、request/replyなどを担う。executorの実行モデルはrclだけで完結せずclient library側にも実装される。

\paragraph{根拠資料}
\begin{itemize}[leftmargin=1.6em]
\item \href{https://docs.ros.org/en/humble/Concepts/About-Internal-Interfaces.html}{ROS 2 Humble公式: Internal ROS 2 interfaces}
\end{itemize}

\subsection{ROS graph、Node、topic、service、action、parameter}
ROS graphは、実行中Nodeと、それらが持つpublisher、subscription、service、actionなどの接続関係である。Pythonクラス、プロセス、ROS graph上のNode名は関連するが同一物ではない。


topicは継続データ向けの非同期一方向publish/subscribe、serviceは短いrequest/response、actionは時間のかかる目標へfeedback、cancel、resultを持たせる。parameterはNodeごとの設定値である。


publisherが送った時点とsubscription callbackが実行される時点は同じとは限らない。通信遅延、QoS queue、executor待ちを挟むため、センサ値にはtimestampとfreshness判定が必要になる。

\paragraph{根拠資料}
\begin{itemize}[leftmargin=1.6em]
\item \href{https://docs.ros.org/en/humble/Tutorials/Beginner-CLI-Tools/Understanding-ROS2-Nodes/Understanding-ROS2-Nodes.html}{ROS 2 Humble公式: Understanding nodes}
\item \href{https://docs.ros.org/en/humble/Concepts/Basic/Interfaces-Topics-Services-Actions.html}{ROS 2 Humble公式: Topics, Services, Actions}
\item \href{https://docs.ros.org/en/humble/Concepts/Basic/About-Parameters.html}{ROS 2 Humble公式: Parameters}
\end{itemize}

\subsection{callback、timer、Executor、spin、Callback Group}
callbackは、メッセージ受信、timer満了、service requestなどのイベントが成立した後でExecutorから呼ばれる関数である。登録時に`self.\_on\_speed`と括弧なしで渡すのは、今実行せず後で呼ぶ関数オブジェクトを渡すためである。


Executorは実行可能になったcallbackを見つけ、Callback Groupの条件を確認して実行する。`spin()`は終了要求まで待機・dispatchを続けるため、通常運転中はmainの次の行へ戻らない。


MutuallyExclusiveCallbackGroupは同じgroup内のcallbackを同時実行させない。別group間はMultiThreadedExecutorで並行実行し得る。groupを分けただけでは共有属性への完全な排他にはならないため、複数groupが同じ状態を読む・書く場合は設計確認が必要である。


\begin{lstlisting}
DDS受信またはtimer満了 -> wait setでready -> Executorが選択 -> Callback Groupが許可 -> worker threadがcallback実行 -> 終了後Executorへ戻る
\end{lstlisting}

\paragraph{根拠資料}
\begin{itemize}[leftmargin=1.6em]
\item \href{https://docs.ros.org/en/rolling/Concepts/Intermediate/About-Executors.html}{ROS 2公式: Executors}
\item \href{https://docs.ros.org/en/humble/Concepts/About-Internal-Interfaces.html}{ROS 2 Humble公式: Internal ROS 2 interfaces}
\end{itemize}

\subsection{rqt\_graph、ros2 CLI、rosbag2をいつ使うか}
rqt\_graphは接続関係を見る道具であり、数値の正しさや更新周期までは保証しない。起動直後、topic名変更後、publisherが複数存在する疑いがある時に使う。


\begin{lstlisting}[language=bash]
ros2 node list
ros2 node info /mpc_node
ros2 topic list -t
ros2 topic info -v /planner/speed_cmd
ros2 topic hz /planner/speed_cmd
ros2 topic echo /planner/status
rqt_graph
\end{lstlisting}


rosbag2はtopicメッセージを時系列のまま記録・再生する。通信不具合、freshness、再計画trigger、実車とSILSの差を再現可能にするため、本番前試験では制御入力だけでなく原因となる全telemetry、status、parameter情報を記録する。


\begin{lstlisting}[language=bash]
ros2 bag record -o outputs/bags/preflight \
  /vehicle/s_km /vehicle/speed_kmh /vehicle/batt_soc \
  /vehicle/batt_temp_c /vehicle/batt_current_a /vehicle/batt_voltage_v \
  /planner/upper_speed_cmd /planner/speed_cmd /planner/status

ros2 bag info outputs/bags/preflight
ros2 bag play outputs/bags/preflight --clock
\end{lstlisting}


bag再生時はQoS互換性、simulation time、外部publisherとの二重入力に注意する。実車Nodeを同時に動かす場合はnamespaceまたはremappingで入力源を明確に分離する。

\paragraph{根拠資料}
\begin{itemize}[leftmargin=1.6em]
\item \href{https://docs.ros.org/en/ros2_packages/humble/api/rqt_graph/index.html}{ROS 2 Humble公式: rqt\_graph}
\item \href{https://docs.ros.org/en/humble/Tutorials/Beginner-CLI-Tools/Recording-And-Playing-Back-Data/Recording-And-Playing-Back-Data.html}{ROS 2 Humble公式: Recording and playing back data}
\item \href{https://docs.ros.org/en/humble/Tutorials/Beginner-CLI-Tools/Understanding-ROS2-Nodes/Understanding-ROS2-Nodes.html}{ROS 2 Humble公式: Understanding nodes}
\end{itemize}

\subsection{天候、route、補間、時刻、単位}
予報は時刻だけの系列か、時刻とroute距離の2次元gridになり得る。現在時刻tと距離sがgrid点の間なら、周囲の値から補間して日射、温度、風を求める。


UTC、現地timezone、naive datetimeを混ぜると数時間のずれが発生する。入力でtimezoneを確定し、内部比較はUTC、表示は現地時刻という役割分担が安全である。


route距離のkm、速度のkm/hとm/s、時間の秒、energyのWhとJを跨ぐため、変換係数3.6、3600、1000の意味を各式で確認する。



\subsection{freshness、filter、guard、fallback、fail-safe}
分散システムでは最後に受け取った値が現在も有効とは限らない。受信時刻とtimeoutからfreshnessを判定し、stale値を計画状態へ無条件に同期しない。


filterはnoiseと一時的な飛び値を抑えるが、遅れを生む。slew limitは指令変化率を制限する。安全guardはsolverのcost罰則とは別に、現在出力へ強制制約を適用する最後の防波堤である。


fallbackは失敗時の代替動作を事前に決める設計である。前回計画保持、物理に基づく決定論的入力、停止、低速制限などから、故障modeごとに選ぶ。fallback発生はstatusとlogへ残し、正常解と区別する。



\subsection{制御、状態、入力、モデル予測制御MPC}
制御対象の内部を表す状態をx、操作入力をu、外乱・予報をwとすると、離散モデルは`x[k+1] = f(x[k], u[k], w[k])`と書ける。ソーラーカーではSoC、電池温度、距離、速度などが状態候補、速度目標や駆動トルクが入力候補になる。


\[
\min_{u_0,\ldots,u_{N-1}} \sum_{k=0}^{N-1}\ell(x_k,u_k,w_k)+V_f(x_N)
\]

MPCは現在状態からNステップ先まで予測し、目的関数と制約を満たす入力系列を求める。ただし実際に適用するのは通常先頭入力だけで、次回は新しい実測状態から再び解く。これがreceding horizonである。


予測モデル、目的関数、制約、ホライズン、solver、初期値のどれかが変わると答えも変わる。「MPCを使う」だけでは仕様は決まらず、これらを単位付きで追う必要がある。

\paragraph{根拠資料}
\begin{itemize}[leftmargin=1.6em]
\item \href{https://docs.scipy.org/doc/scipy/reference/generated/scipy.optimize.minimize.html}{SciPy公式: scipy.optimize.minimize}
\end{itemize}



        \section{関数・クラスを上から順に解説}
        \subsection{L16 関数 \_timestamp\_ns}
            \begin{itemize}[leftmargin=1.6em]
              \item 行範囲: L16--L17
              \item 所属: module直下
              \item 定義: \path|_timestamp_ns(value) -> int|
              \item docstring: 明示されていない。
              \item このブロックが直接呼ぶ主な関数・メソッド: Timestamp, int
              \item 戻り値の要点: int(pd.Timestamp(value).value)
              \item 主なローカル代入名: 特になし。
              \item 読み取る主なインスタンス属性: 特になし。
              \item 更新する主なインスタンス属性: 特になし。
              \item 明示的に送出する例外: 明示的raiseはない。
              \item 制御構造の規模: 条件分岐 0、ループ 0、try 0
            \end{itemize}
            \paragraph{この定義を読むためのPython構文}
            \begin{itemize}[leftmargin=1.6em]
            \item `->`以後は戻り値の型注釈であり、実行時の値を自動変換する命令ではない。
            \end{itemize}
            \paragraph{上から順の処理}
            \begin{enumerate}[leftmargin=1.8em]
            \item int(pd.Timestamp(value).value) を返す。
            \end{enumerate}
            \paragraph{代表コード断片}
            \begin{lstlisting}[language=Python]
def _timestamp_ns(value) -> int:
    return int(pd.Timestamp(value).value)
\end{lstlisting}

\subsection{L20 クラス WindCorrectionNode}
            \begin{itemize}[leftmargin=1.6em]
              \item 行範囲: L20--L151
              \item 所属: module直下
              \item 定義: \path|WindCorrectionNode(bases=Node)|
              \item docstring: 明示されていない。
              \item このブロックが直接呼ぶ主な関数・メソッド: 特に明示されていない。
              \item 戻り値の要点: 戻り値が無いか、return 文が明示されていない。
              \item 主なローカル代入名: 特になし。
              \item 読み取る主なインスタンス属性: 特になし。
              \item 更新する主なインスタンス属性: 特になし。
              \item 明示的に送出する例外: 明示的raiseはない。
              \item 制御構造の規模: 条件分岐 0、ループ 0、try 0
            \end{itemize}
            \paragraph{この定義を読むためのPython構文}
            \begin{itemize}[leftmargin=1.6em]
            \item class文は新しい型を定義する。丸括弧内は継承する基底クラスである。
\item 内包表記は、反復・条件・値生成を一つの式で記述してcollectionを作る。
\item try文は例外が起き得る処理と、異常時または終了時の経路を分ける。
            \end{itemize}
            \paragraph{上から順の処理}
            \begin{enumerate}[leftmargin=1.8em]
            \item 関数 \_\_init\_\_ を定義する。
\item 関数 \_on\_headwind を定義する。
\item 関数 \_on\_s を定義する。
\item 関数 \_load\_base を定義する。
\item 関数 \_interp\_headwind\_now を定義する。
\item 関数 \_step を定義する。
            \end{enumerate}
            \paragraph{代表コード断片}
            \begin{lstlisting}[language=Python]
class WindCorrectionNode(Node):
    def __init__(self) -> None:
        super().__init__("wind_correction_node")
        self.declare_parameter("forecast_csv", "")
        self.declare_parameter("corrected_forecast_csv", "outputs/runtime/live_forecast_corrected.csv")
        self.declare_parameter("forecast_time_tz", "Australia/Darwin")
        self.declare_parameter("publish_period_sec", 30.0)
        self.declare_parameter("measurement_sigma_ms", 1.0)
        self.declare_parameter("correlation_distance_km", 300.0)
        self.declare_parameter("fallback_correlation_time_h", 3.0)
        self.declare_parameter("forecast_sigma0_ms", 1.5)
        self.declare_parameter("forecast_variance_growth_per_hour", 0.05)
        self.declare_parameter("planning_quantile", 0.5)
        self.declare_parameter("confidence_z", 1.96)
        self.declare_parameter("min_sigma_ms", 0.2)
        self.declare_parameter("preferred_source", "auto")
        self.declare_parameter("use_exp_distance_decay", True)

        self.forecast_csv = Path(str(self.get_parameter("forecast_csv").value or "")).expanduser()
        self.corrected_csv = Path(str(self.get_parameter("corrected_forecast_csv").value or "")).expanduser()
        self.forecast_time_tz = str(self.get_parameter("forecast_time_tz").value or "Australia/Darwin")
        self.measurement_sigma = max(0.0, float(self.get_parameter("measurement_sigma_ms").value))
        self.correlation_distance_km = max(1.0, float(self.get_parameter("correlation_distance_km").value))
        self.correlation_time_h = max(0.25, float(self.get_parameter("fallback_correlation_time_h").value))
        self.forecast_sigma0 = max(0.0, float(self.get_parameter("forecast_sigma0_ms").value))
        self.variance_growth = max(0.0, float(self.get_parameter("forecast_variance_growth_per_hour").value))
        self.planning_quantile = min(0.999, max(0.001, float(self.get_parameter("planning_quantile").value)))
        self.confidence_z = max(0.1, float(self.get_parameter("confidence_z").value))
        self.min_sigma = max(0.0, float(self.get_parameter("min_sigma_ms").value))
        self.use_exp_distance_decay = bool(self.get_parameter("use_exp_distance_decay").value)
        self.quantile_z = float(NormalDist().inv_cdf(self.planning_quantile))

        self.latest_headwind = math.nan
        self.latest_s_km = math.nan
        self.latest_bias = 0.0
...
\end{lstlisting}

\subsection{L21 関数 WindCorrectionNode.\_\_init\_\_}
            \begin{itemize}[leftmargin=1.6em]
              \item 行範囲: L21--L62
              \item 所属: \path|WindCorrectionNode|
              \item 定義: \path|__init__(self) -> None|
              \item docstring: 明示されていない。
              \item このブロックが直接呼ぶ主な関数・メソッド: NormalDist, Path, \_\_init\_\_, bool, create\_publisher, create\_subscription, create\_timer, declare\_parameter, expanduser, float, get\_logger, get\_parameter
              \item 戻り値の要点: 戻り値が無いか、return 文が明示されていない。
              \item 主なローカル代入名: 特になし。
              \item 読み取る主なインスタンス属性: self.\_on\_headwind, self.\_on\_s, self.\_step, self.corrected\_csv, self.create\_publisher, self.create\_subscription, self.create\_timer, self.declare\_parameter, self.forecast\_csv, self.get\_logger, self.get\_parameter, self.planning\_quantile
              \item 更新する主なインスタンス属性: self.base\_df, self.base\_mtime, self.confidence\_z, self.corrected\_csv, self.correlation\_distance\_km, self.correlation\_time\_h, self.forecast\_csv, self.forecast\_sigma0, self.forecast\_time\_tz, self.latest\_bias, self.latest\_headwind, self.latest\_s\_km, self.measurement\_sigma, self.min\_sigma, self.planning\_quantile, self.pub\_headwind, self.pub\_status, self.quantile\_z, self.timer, self.use\_exp\_distance\_decay, self.variance\_growth
              \item 明示的に送出する例外: 明示的raiseはない。
              \item 制御構造の規模: 条件分岐 0、ループ 0、try 0
            \end{itemize}
            \paragraph{この定義を読むためのPython構文}
            \begin{itemize}[leftmargin=1.6em]
            \item 第1引数selfは、このメソッドを呼んだインスタンス自身を受け取る慣習名である。
\item `->`以後は戻り値の型注釈であり、実行時の値を自動変換する命令ではない。
            \end{itemize}
            \paragraph{上から順の処理}
            \begin{enumerate}[leftmargin=1.8em]
            \item super().\_\_init\_\_(...) を実行する。
\item self.declare\_parameter(...) を実行する。
\item self.declare\_parameter(...) を実行する。
\item self.declare\_parameter(...) を実行する。
\item self.declare\_parameter(...) を実行する。
\item self.declare\_parameter(...) を実行する。
\item self.declare\_parameter(...) を実行する。
\item self.declare\_parameter(...) を実行する。
\item self.declare\_parameter(...) を実行する。
\item self.declare\_parameter(...) を実行する。
\item self.declare\_parameter(...) を実行する。
\item self.declare\_parameter(...) を実行する。
\item self.declare\_parameter(...) を実行する。
\item self.declare\_parameter(...) を実行する。
\item self.declare\_parameter(...) を実行する。
\item self.forecast\_csv に Path(str(self.get\_parameter('forecast\_csv').value or '')).expanduser() の結果を代入する。
\item self.corrected\_csv に Path(str(self.get\_parameter('corrected\_forecast\_csv').value or '')).expanduser() の結果を代入する。
\item self.forecast\_time\_tz に str(self.get\_parameter('forecast\_time\_tz').value or 'Australia/Darwin') の結果を代入する。
\item self.measurement\_sigma に max(0.0, float(self.get\_parameter('measurement\_sigma\_ms').value)) の結果を代入する。
\item self.correlation\_distance\_km に max(1.0, float(self.get\_parameter('correlation\_distance\_km').value)) の結果を代入する。
\item self.correlation\_time\_h に max(0.25, float(self.get\_parameter('fallback\_correlation\_time\_h').value)) の結果を代入する。
\item self.forecast\_sigma0 に max(0.0, float(self.get\_parameter('forecast\_sigma0\_ms').value)) の結果を代入する。
\item self.variance\_growth に max(0.0, float(self.get\_parameter('forecast\_variance\_growth\_per\_hour').value)) の結果を代入する。
\item self.planning\_quantile に min(0.999, max(0.001, float(self.get\_parameter('planning\_quantile').value))) の結果を代入する。
\item self.confidence\_z に max(0.1, float(self.get\_parameter('confidence\_z').value)) の結果を代入する。
\item self.min\_sigma に max(0.0, float(self.get\_parameter('min\_sigma\_ms').value)) の結果を代入する。
\item self.use\_exp\_distance\_decay に bool(self.get\_parameter('use\_exp\_distance\_decay').value) の結果を代入する。
\item self.quantile\_z に float(NormalDist().inv\_cdf(self.planning\_quantile)) の結果を代入する。
\item self.latest\_headwind に math.nan の結果を代入する。
\item self.latest\_s\_km に math.nan の結果を代入する。
\item self.latest\_bias に 0.0 の結果を代入する。
\item self.base\_df に None の結果を代入する。
\item self.base\_mtime に None の結果を代入する。
\item self.pub\_headwind に self.create\_publisher(Float32, '/weather/headwind\_corrected\_ms', 10) の結果を代入する。
\item self.pub\_status に self.create\_publisher(String, '/weather/wind\_correction\_status', 10) の結果を代入する。
\item self.create\_subscription(...) を実行する。
\item self.create\_subscription(...) を実行する。
\item self.timer に self.create\_timer(max(1.0, float(self.get\_parameter('publish\_period\_sec').value)), self.\_step) の結果を代入する。
\item self.get\_logger().info(...) を実行する。
            \end{enumerate}
            \paragraph{代表コード断片}
            \begin{lstlisting}[language=Python]
    def __init__(self) -> None:
        super().__init__("wind_correction_node")
        self.declare_parameter("forecast_csv", "")
        self.declare_parameter("corrected_forecast_csv", "outputs/runtime/live_forecast_corrected.csv")
        self.declare_parameter("forecast_time_tz", "Australia/Darwin")
        self.declare_parameter("publish_period_sec", 30.0)
        self.declare_parameter("measurement_sigma_ms", 1.0)
        self.declare_parameter("correlation_distance_km", 300.0)
        self.declare_parameter("fallback_correlation_time_h", 3.0)
        self.declare_parameter("forecast_sigma0_ms", 1.5)
        self.declare_parameter("forecast_variance_growth_per_hour", 0.05)
        self.declare_parameter("planning_quantile", 0.5)
        self.declare_parameter("confidence_z", 1.96)
        self.declare_parameter("min_sigma_ms", 0.2)
        self.declare_parameter("preferred_source", "auto")
        self.declare_parameter("use_exp_distance_decay", True)

        self.forecast_csv = Path(str(self.get_parameter("forecast_csv").value or "")).expanduser()
        self.corrected_csv = Path(str(self.get_parameter("corrected_forecast_csv").value or "")).expanduser()
        self.forecast_time_tz = str(self.get_parameter("forecast_time_tz").value or "Australia/Darwin")
        self.measurement_sigma = max(0.0, float(self.get_parameter("measurement_sigma_ms").value))
        self.correlation_distance_km = max(1.0, float(self.get_parameter("correlation_distance_km").value))
        self.correlation_time_h = max(0.25, float(self.get_parameter("fallback_correlation_time_h").value))
        self.forecast_sigma0 = max(0.0, float(self.get_parameter("forecast_sigma0_ms").value))
        self.variance_growth = max(0.0, float(self.get_parameter("forecast_variance_growth_per_hour").value))
        self.planning_quantile = min(0.999, max(0.001, float(self.get_parameter("planning_quantile").value)))
        self.confidence_z = max(0.1, float(self.get_parameter("confidence_z").value))
        self.min_sigma = max(0.0, float(self.get_parameter("min_sigma_ms").value))
        self.use_exp_distance_decay = bool(self.get_parameter("use_exp_distance_decay").value)
        self.quantile_z = float(NormalDist().inv_cdf(self.planning_quantile))

        self.latest_headwind = math.nan
        self.latest_s_km = math.nan
        self.latest_bias = 0.0
        self.base_df = None
...
\end{lstlisting}

\subsection{L64 関数 WindCorrectionNode.\_on\_headwind}
            \begin{itemize}[leftmargin=1.6em]
              \item 行範囲: L64--L65
              \item 所属: \path|WindCorrectionNode|
              \item 定義: \path|_on_headwind(self, msg: Float32) -> None|
              \item docstring: 明示されていない。
              \item このブロックが直接呼ぶ主な関数・メソッド: float
              \item 戻り値の要点: 戻り値が無いか、return 文が明示されていない。
              \item 主なローカル代入名: 特になし。
              \item 読み取る主なインスタンス属性: 特になし。
              \item 更新する主なインスタンス属性: self.latest\_headwind
              \item 明示的に送出する例外: 明示的raiseはない。
              \item 制御構造の規模: 条件分岐 0、ループ 0、try 0
            \end{itemize}
            \paragraph{この定義を読むためのPython構文}
            \begin{itemize}[leftmargin=1.6em]
            \item 第1引数selfは、このメソッドを呼んだインスタンス自身を受け取る慣習名である。
\item `->`以後は戻り値の型注釈であり、実行時の値を自動変換する命令ではない。
            \end{itemize}
            \paragraph{上から順の処理}
            \begin{enumerate}[leftmargin=1.8em]
            \item self.latest\_headwind に float(msg.data) の結果を代入する。
            \end{enumerate}
            \paragraph{代表コード断片}
            \begin{lstlisting}[language=Python]
    def _on_headwind(self, msg: Float32) -> None:
        self.latest_headwind = float(msg.data)
\end{lstlisting}

\subsection{L67 関数 WindCorrectionNode.\_on\_s}
            \begin{itemize}[leftmargin=1.6em]
              \item 行範囲: L67--L68
              \item 所属: \path|WindCorrectionNode|
              \item 定義: \path|_on_s(self, msg: Float32) -> None|
              \item docstring: 明示されていない。
              \item このブロックが直接呼ぶ主な関数・メソッド: float
              \item 戻り値の要点: 戻り値が無いか、return 文が明示されていない。
              \item 主なローカル代入名: 特になし。
              \item 読み取る主なインスタンス属性: 特になし。
              \item 更新する主なインスタンス属性: self.latest\_s\_km
              \item 明示的に送出する例外: 明示的raiseはない。
              \item 制御構造の規模: 条件分岐 0、ループ 0、try 0
            \end{itemize}
            \paragraph{この定義を読むためのPython構文}
            \begin{itemize}[leftmargin=1.6em]
            \item 第1引数selfは、このメソッドを呼んだインスタンス自身を受け取る慣習名である。
\item `->`以後は戻り値の型注釈であり、実行時の値を自動変換する命令ではない。
            \end{itemize}
            \paragraph{上から順の処理}
            \begin{enumerate}[leftmargin=1.8em]
            \item self.latest\_s\_km に float(msg.data) の結果を代入する。
            \end{enumerate}
            \paragraph{代表コード断片}
            \begin{lstlisting}[language=Python]
    def _on_s(self, msg: Float32) -> None:
        self.latest_s_km = float(msg.data)
\end{lstlisting}

\subsection{L70 関数 WindCorrectionNode.\_load\_base}
            \begin{itemize}[leftmargin=1.6em]
              \item 行範囲: L70--L89
              \item 所属: \path|WindCorrectionNode|
              \item 定義: \path|_load_base(self) -> None|
              \item docstring: 明示されていない。
              \item このブロックが直接呼ぶ主な関数・メソッド: getattr, getmtime, read\_csv, to\_datetime, tz\_convert, tz\_localize
              \item 戻り値の要点: 戻り値が無いか、return 文が明示されていない。
              \item 主なローカル代入名: df, mtime, t
              \item 読み取る主なインスタンス属性: self.base\_df, self.base\_mtime, self.forecast\_csv, self.forecast\_time\_tz
              \item 更新する主なインスタンス属性: self.base\_df, self.base\_mtime
              \item 明示的に送出する例外: 明示的raiseはない。
              \item 制御構造の規模: 条件分岐 3、ループ 0、try 2
            \end{itemize}
            \paragraph{この定義を読むためのPython構文}
            \begin{itemize}[leftmargin=1.6em]
            \item 第1引数selfは、このメソッドを呼んだインスタンス自身を受け取る慣習名である。
\item `->`以後は戻り値の型注釈であり、実行時の値を自動変換する命令ではない。
\item try文は例外が起き得る処理と、異常時または終了時の経路を分ける。
            \end{itemize}
            \paragraph{上から順の処理}
            \begin{enumerate}[leftmargin=1.8em]
            \item 例外処理を伴う try ブロックを実行する。
\item   mtime に os.path.getmtime(self.forecast\_csv) の結果を代入する。
\item   Exceptionを捕捉した場合:
\item    を返す。
\item 条件 self.base\_mtime is not None and mtime <= self.base\_mtime and (self.base\_df is not None) を判定し、真なら内部処理を行う。
\item    を返す。
\item df に pd.read\_csv(self.forecast\_csv) の結果を代入する。
\item 条件 'time' in df.columns を判定し、真なら内部処理を行う。
\item   t に pd.to\_datetime(df['time'], format='mixed', errors='coerce') の結果を代入する。
\item   条件 getattr(t.dt, 'tz', None) is None を判定し、真なら内部処理を行う。
\item     例外処理を伴う try ブロックを実行する。
\item       t に t.dt.tz\_localize(self.forecast\_time\_tz, ambiguous='NaT', nonexistent='NaT').dt.tz\_convert('UTC') の結果を代入する。
\item       Exceptionを捕捉した場合:
\item       t に t.dt.tz\_localize('UTC') の結果を代入する。
\item     上の条件が偽の場合:
\item     t に t.dt.tz\_convert('UTC') の結果を代入する。
\item   df['time'] に t の結果を代入する。
\item self.base\_df に df の結果を代入する。
\item self.base\_mtime に mtime の結果を代入する。
            \end{enumerate}
            \paragraph{代表コード断片}
            \begin{lstlisting}[language=Python]
    def _load_base(self) -> None:
        try:
            mtime = os.path.getmtime(self.forecast_csv)
        except Exception:
            return
        if self.base_mtime is not None and mtime <= self.base_mtime and self.base_df is not None:
            return
        df = pd.read_csv(self.forecast_csv)
        if "time" in df.columns:
            t = pd.to_datetime(df["time"], format="mixed", errors="coerce")
            if getattr(t.dt, "tz", None) is None:
                try:
                    t = t.dt.tz_localize(self.forecast_time_tz, ambiguous="NaT", nonexistent="NaT").dt.tz_convert("UTC")
                except Exception:
                    t = t.dt.tz_localize("UTC")
            else:
                t = t.dt.tz_convert("UTC")
            df["time"] = t
        self.base_df = df
        self.base_mtime = mtime
\end{lstlisting}

\subsection{L91 関数 WindCorrectionNode.\_interp\_headwind\_now}
            \begin{itemize}[leftmargin=1.6em]
              \item 行範囲: L91--L103
              \item 所属: \path|WindCorrectionNode|
              \item 定義: \path|_interp_headwind_now(self, now_utc: datetime) -> float|
              \item docstring: 明示されていない。
              \item このブロックが直接呼ぶ主な関数・メソッド: Timestamp, any, float, get, int, len, max, min, notna, searchsorted
              \item 戻り値の要点: 0.0 / float(row.get('headwind\_ms', 0.0)) / 0.0
              \item 主なローカル代入名: idx, row, work
              \item 読み取る主なインスタンス属性: self.base\_df
              \item 更新する主なインスタンス属性: 特になし。
              \item 明示的に送出する例外: 明示的raiseはない。
              \item 制御構造の規模: 条件分岐 2、ループ 0、try 1
            \end{itemize}
            \paragraph{この定義を読むためのPython構文}
            \begin{itemize}[leftmargin=1.6em]
            \item 第1引数selfは、このメソッドを呼んだインスタンス自身を受け取る慣習名である。
\item `->`以後は戻り値の型注釈であり、実行時の値を自動変換する命令ではない。
\item try文は例外が起き得る処理と、異常時または終了時の経路を分ける。
            \end{itemize}
            \paragraph{上から順の処理}
            \begin{enumerate}[leftmargin=1.8em]
            \item 条件 self.base\_df is None or self.base\_df.empty を判定し、真なら内部処理を行う。
\item   0.0 を返す。
\item work に self.base\_df の結果を代入する。
\item 条件 'time' in work.columns and work['time'].notna().any() を判定し、真なら内部処理を行う。
\item   idx に int(max(0, min(len(work) - 1, work['time'].searchsorted(pd.Timestamp(now\_utc), side='right') - 1))) の結果を代入する。
\item   row に work.iloc[idx] の結果を代入する。
\item   上の条件が偽の場合:
\item   row に work.iloc[0] の結果を代入する。
\item 例外処理を伴う try ブロックを実行する。
\item   float(row.get('headwind\_ms', 0.0)) を返す。
\item   Exceptionを捕捉した場合:
\item   0.0 を返す。
            \end{enumerate}
            \paragraph{代表コード断片}
            \begin{lstlisting}[language=Python]
    def _interp_headwind_now(self, now_utc: datetime) -> float:
        if self.base_df is None or self.base_df.empty:
            return 0.0
        work = self.base_df
        if "time" in work.columns and work["time"].notna().any():
            idx = int(max(0, min(len(work) - 1, work["time"].searchsorted(pd.Timestamp(now_utc), side="right") - 1)))
            row = work.iloc[idx]
        else:
            row = work.iloc[0]
        try:
            return float(row.get("headwind_ms", 0.0))
        except Exception:
            return 0.0
\end{lstlisting}

\subsection{L105 関数 WindCorrectionNode.\_step}
            \begin{itemize}[leftmargin=1.6em]
              \item 行範囲: L105--L151
              \item 所属: \path|WindCorrectionNode|
              \item 定義: \path|_step(self) -> None|
              \item docstring: 明示されていない。
              \item このブロックが直接呼ぶ主な関数・メソッド: Float32, String, Timedelta, Timestamp, \_interp\_headwind\_now, \_load\_base, abs, append, astimezone, copy, exp, fillna
              \item 戻り値の要点: 戻り値が無いか、return 文が明示されていない。
              \item 主なローカル代入名: base, base\_now, corrected, corrected\_now, corrected\_values, current\_s, decay, df, dist\_decay, ds, dt\_h, i, innovation, mean, now\_utc, out, row, sigma, sigma\_values, t\_utc
              \item 読み取る主なインスタンス属性: self.\_interp\_headwind\_now, self.\_load\_base, self.base\_df, self.confidence\_z, self.corrected\_csv, self.correlation\_distance\_km, self.correlation\_time\_h, self.forecast\_sigma0, self.forecast\_time\_tz, self.latest\_bias, self.latest\_headwind, self.latest\_s\_km, self.measurement\_sigma, self.min\_sigma, self.pub\_headwind, self.pub\_status, self.quantile\_z, self.use\_exp\_distance\_decay, self.variance\_growth
              \item 更新する主なインスタンス属性: self.latest\_bias
              \item 明示的に送出する例外: 明示的raiseはない。
              \item 制御構造の規模: 条件分岐 5、ループ 1、try 0
            \end{itemize}
            \paragraph{この定義を読むためのPython構文}
            \begin{itemize}[leftmargin=1.6em]
            \item 第1引数selfは、このメソッドを呼んだインスタンス自身を受け取る慣習名である。
\item `->`以後は戻り値の型注釈であり、実行時の値を自動変換する命令ではない。
\item 内包表記は、反復・条件・値生成を一つの式で記述してcollectionを作る。
            \end{itemize}
            \paragraph{上から順の処理}
            \begin{enumerate}[leftmargin=1.8em]
            \item self.\_load\_base(...) を実行する。
\item 条件 self.base\_df is None or self.base\_df.empty を判定し、真なら内部処理を行う。
\item    を返す。
\item now\_utc に datetime.now(timezone.utc) の結果を代入する。
\item base\_now に self.\_interp\_headwind\_now(now\_utc) の結果を代入する。
\item 条件 math.isfinite(self.latest\_headwind) を判定し、真なら内部処理を行う。
\item   innovation に float(self.latest\_headwind - base\_now) の結果を代入する。
\item   self.latest\_bias に 0.7 * float(self.latest\_bias) + 0.3 * innovation の結果を代入する。
\item corrected\_now に base\_now + self.latest\_bias の結果を代入する。
\item self.pub\_headwind.publish(...) を実行する。
\item df に self.base\_df.copy() の結果を代入する。
\item 条件 'time' not in df.columns を判定し、真なら内部処理を行う。
\item   df['time'] に [now\_utc + pd.Timedelta(minutes=10 * i) for i in range(len(df))] の結果を代入する。
\item 条件 's\_km' not in df.columns を判定し、真なら内部処理を行う。
\item   df['s\_km'] に float(self.latest\_s\_km) if math.isfinite(self.latest\_s\_km) else 0.0 の結果を代入する。
\item current\_s に float(self.latest\_s\_km) if math.isfinite(self.latest\_s\_km) else float(df['s\_km'].iloc[0]) の結果を代入する。
\item corrected\_values に [] の結果を代入する。
\item sigma\_values に [] の結果を代入する。
\item df.itertuples(index=False) を順に走査し、各要素を row に入れて処理する。
\item   t\_utc に getattr(row, 'time') の結果を代入する。
\item   dt\_h に max(0.0, (pd.Timestamp(t\_utc).to\_pydatetime().astimezone(timezone.utc) - now\_utc).total\_seconds() / 3600.0) の結果を代入する。
\item   ds に abs(float(getattr(row, 's\_km', current\_s)) - current\_s) の結果を代入する。
\item   time\_decay に math.exp(-dt\_h / self.correlation\_time\_h) の結果を代入する。
\item   dist\_decay に math.exp(-ds / self.correlation\_distance\_km) if self.use\_exp\_distance\_decay else 1.0 / (1.0 + ds / self.correlation\_distance\_km) の結果を代入する。
\item   decay に time\_decay * dist\_decay の結果を代入する。
\item   base に float(getattr(row, 'headwind\_ms', 0.0)) の結果を代入する。
\item   mean に base + self.latest\_bias * decay の結果を代入する。
\item   sigma に max(self.min\_sigma, math.sqrt(self.forecast\_sigma0 ** 2 + self.variance\_growth * dt\_h) + self.measurement\_sigma * (1.0 - decay)) の結果を代入する。
\item   corrected に mean + self.quantile\_z * sigma の結果を代入する。
\item   corrected\_values.append(...) を実行する。
\item   sigma\_values.append(...) を実行する。
\item df['headwind\_base\_ms'] に pd.to\_numeric(df.get('headwind\_ms', 0.0), errors='coerce').fillna(0.0) の結果を代入する。
\item df['headwind\_mean\_ms'] に corrected\_values の結果を代入する。
\item df['headwind\_sigma\_ms'] に sigma\_values の結果を代入する。
\item df['headwind\_lower\_ms'] に df['headwind\_mean\_ms'] - self.confidence\_z * df['headwind\_sigma\_ms'] の結果を代入する。
\item df['headwind\_upper\_ms'] に df['headwind\_mean\_ms'] + self.confidence\_z * df['headwind\_sigma\_ms'] の結果を代入する。
\item df['headwind\_ms'] に corrected\_values の結果を代入する。
\item out に df.copy() の結果を代入する。
\item 条件 pd.api.types.is\_datetime64\_any\_dtype(out['time']) を判定し、真なら内部処理を行う。
\item   out['time'] に out['time'].dt.tz\_convert(self.forecast\_time\_tz).dt.strftime('\%Y-\%m-\%d \%H:\%M:\%S') の結果を代入する。
\item self.corrected\_csv.parent.mkdir(...) を実行する。
\item out.to\_csv(...) を実行する。
\item self.pub\_status.publish(...) を実行する。
            \end{enumerate}
            \paragraph{代表コード断片}
            \begin{lstlisting}[language=Python]
    def _step(self) -> None:
        self._load_base()
        if self.base_df is None or self.base_df.empty:
            return
        now_utc = datetime.now(timezone.utc)
        base_now = self._interp_headwind_now(now_utc)
        if math.isfinite(self.latest_headwind):
            innovation = float(self.latest_headwind - base_now)
            self.latest_bias = 0.7 * float(self.latest_bias) + 0.3 * innovation
        corrected_now = base_now + self.latest_bias
        self.pub_headwind.publish(Float32(data=float(corrected_now)))

        df = self.base_df.copy()
        if "time" not in df.columns:
            df["time"] = [now_utc + pd.Timedelta(minutes=10 * i) for i in range(len(df))]
        if "s_km" not in df.columns:
            df["s_km"] = float(self.latest_s_km) if math.isfinite(self.latest_s_km) else 0.0
        current_s = float(self.latest_s_km) if math.isfinite(self.latest_s_km) else float(df["s_km"].iloc[0])

        corrected_values = []
        sigma_values = []
        for row in df.itertuples(index=False):
            t_utc = getattr(row, "time")
            dt_h = max(0.0, (pd.Timestamp(t_utc).to_pydatetime().astimezone(timezone.utc) - now_utc).total_seconds() / 3600.0)
            ds = abs(float(getattr(row, "s_km", current_s)) - current_s)
            time_decay = math.exp(-dt_h / self.correlation_time_h)
            dist_decay = math.exp(-ds / self.correlation_distance_km) if self.use_exp_distance_decay else 1.0 / (1.0 + ds / self.correlation_distance_km)
            decay = time_decay * dist_decay
            base = float(getattr(row, "headwind_ms", 0.0))
            mean = base + self.latest_bias * decay
            sigma = max(self.min_sigma, math.sqrt(self.forecast_sigma0 ** 2 + self.variance_growth * dt_h) + self.measurement_sigma * (1.0 - decay))
            corrected = mean + self.quantile_z * sigma
            corrected_values.append(corrected)
            sigma_values.append(sigma)

...
\end{lstlisting}

\subsection{L154 関数 main}
            \begin{itemize}[leftmargin=1.6em]
              \item 行範囲: L154--L159
              \item 所属: module直下
              \item 定義: \path|main() -> None|
              \item docstring: 明示されていない。
              \item このブロックが直接呼ぶ主な関数・メソッド: WindCorrectionNode, destroy\_node, init, shutdown, spin
              \item 戻り値の要点: 戻り値が無いか、return 文が明示されていない。
              \item 主なローカル代入名: node
              \item 読み取る主なインスタンス属性: 特になし。
              \item 更新する主なインスタンス属性: 特になし。
              \item 明示的に送出する例外: 明示的raiseはない。
              \item 制御構造の規模: 条件分岐 0、ループ 0、try 0
            \end{itemize}
            \paragraph{この定義を読むためのPython構文}
            \begin{itemize}[leftmargin=1.6em]
            \item `->`以後は戻り値の型注釈であり、実行時の値を自動変換する命令ではない。
            \end{itemize}
            \paragraph{上から順の処理}
            \begin{enumerate}[leftmargin=1.8em]
            \item rclpy.init(...) を実行する。
\item node に WindCorrectionNode() の結果を代入する。
\item rclpy.spin(...) を実行する。
\item node.destroy\_node(...) を実行する。
\item rclpy.shutdown(...) を実行する。
            \end{enumerate}
            \paragraph{代表コード断片}
            \begin{lstlisting}[language=Python]
def main() -> None:
    rclpy.init()
    node = WindCorrectionNode()
    rclpy.spin(node)
    node.destroy_node()
    rclpy.shutdown()
\end{lstlisting}

        \subsection{設定値と入力パラメータ}
            \begin{longtable}{p{0.07\linewidth} p{0.35\linewidth} p{0.45\linewidth}}
            \toprule
            行 & 名前 & 既定値・補足 \\
            \midrule
            23 & \path|forecast_csv| &  \\
24 & \path|corrected_forecast_csv| & outputs/runtime/live\_forecast\_corrected.csv \\
25 & \path|forecast_time_tz| & Australia/Darwin \\
26 & \path|publish_period_sec| & 30.0 \\
27 & \path|measurement_sigma_ms| & 1.0 \\
28 & \path|correlation_distance_km| & 300.0 \\
29 & \path|fallback_correlation_time_h| & 3.0 \\
30 & \path|forecast_sigma0_ms| & 1.5 \\
31 & \path|forecast_variance_growth_per_hour| & 0.05 \\
32 & \path|planning_quantile| & 0.5 \\
33 & \path|confidence_z| & 1.96 \\
34 & \path|min_sigma_ms| & 0.2 \\
35 & \path|preferred_source| & auto \\
36 & \path|use_exp_distance_decay| & True \\
            \bottomrule
            \end{longtable}

        \subsection{ROS topic I/O}
            \begin{longtable}{p{0.07\linewidth} p{0.18\linewidth} p{0.34\linewidth} p{0.28\linewidth}}
            \toprule
            行 & 種別 & topic & callback / 補足 \\
            \midrule
            57 & publisher & \path|/weather/headwind_corrected_ms| & publisher \\
58 & publisher & \path|/weather/wind_correction_status| & publisher \\
59 & subscription & \path|/weather/headwind_meas_ms| & self.\_on\_headwind \\
60 & subscription & \path|/vehicle/s_km| & self.\_on\_s \\
            \bottomrule
            \end{longtable}







        \section{処理の流れ}
        \begin{enumerate}[leftmargin=1.7em]
        \item 初期化時に設定値や入力パスを読み込む。
\item publisher / subscription / timer を準備する。
\item timer callback 周期で主処理を進める。
        \end{enumerate}

        \section{このファイルの位置づけ}
        この資料群では、\path|mpc_solarcar/wind_correction_node.py| を
        \textbf{runtime node}の中核として扱う。
        したがって、ここで扱う処理は単独で完結せず、
        前段の profile / launch / telemetry と後段の planner / logger / report へ繋がる。

        \end{document}
